\section{Few-shot 阈值、Inverse Scaling 与 U-shaped Scaling}

有了操作性定义，下一步不是立即寻找神秘机制，而是学会读图。Few-shot prompting 把若干 input--output exemplar 放入上下文，让 frozen language model 继续预测答案；如果小模型只会猜、大模型能利用 exemplar 完成任务，便构成课堂所说的 few-shot emergence。

\subsection{Few-shot prompting 如何变成可测实验}

实验需要明确 prompt、未见 input、target label 和 chance baseline。若二分类任务的随机准确率为 $0.5$，小模型在 $0.5$ 左右并不证明它“一点也不懂”，但至少说明当前 interface 没有把知识转化为可用性能；大模型持续超过 baseline 才说明这套 interface 下的能力可测。换言之，能力判断不是读一段模型输出后的主观印象，而是把同一种 elicitation procedure 放到一系列 scale checkpoints 上比较，并检查改进是否超过抽样噪声、类别不平衡和偶然 prompt 选择能够解释的范围。

\lecturefigure{slide-07.jpg}{Few-shot prompted task 的涌现定义}{官方课件第 7 页；Wei et al., 2022}

对 $K$ 类平衡分类，可用 chance-normalized gain 辅助比较：
\[
G_{\mathrm{chance}}=\frac{a-1/K}{1-1/K},
\]
其中 $a$ 是准确率，$1/K$ 是随机基线。$G_{\mathrm{chance}}=0$ 表示没有超过随机，$G_{\mathrm{chance}}=1$ 表示满分。它便于跨不同类别数的任务观察“超过 chance 的程度”，但仍不消除 prompt sensitivity 与样本方差。

\lecturefigure{slide-08.jpg}{多个任务上的 few-shot emergence 曲线}{官方课件第 8 页；Wei et al., 2022}

\begin{knowledgebox}{读图：先看 baseline，再看拐点}
每个小图的横轴是训练 FLOPs 或模型规模，纵轴是某个任务 metric。第一步找 random baseline；第二步看小模型点是否围绕 baseline 波动；第三步判断大模型是否连续多个点超过 baseline，而不是只挑一个噪声点；第四步比较不同模型家族的拐点是否一致。所谓“突然”，描述的是稀疏采样下的任务曲线，不代表内部表示在相邻两个 checkpoint 之间瞬间生成。
\end{knowledgebox}

\teachervoice{讲者逐轴解释这些图：横轴是规模，纵轴是性能，虚线或已知数值是随机水平。课堂强调的不是“肉眼看到弯曲”，而是小模型未超过随机、大模型开始稳定超过随机的判据。}

\subsection{两个具体任务：知识测验与语音转写}

抽象小图容易让读者忽略 task mechanics，因此课件接着给出两个更具体的例子。一个是多学科选择题，另一个是把英语句子转为 International Phonetic Alphabet（IPA，国际音标）；二者都要求模型把上下文示例转成新的输出格式。

\lecturefigure{slide-09.jpg}{多任务选择题中的 few-shot 阈值现象}{官方课件第 9 页；Hendrycks et al., 2020 数据}

这类 benchmark 的 raw accuracy 混合了知识、读题、选项比较和输出格式。若小模型无法稳定执行其中任一环节，总分会贴近随机；规模提升可能让多个环节同时变得“足够可靠”，从而在 aggregate metric 上形成明显拐点。图不能单独告诉我们究竟是哪一环先改善。

\lecturefigure{slide-10.jpg}{英语到 IPA 转写任务中的涌现表现}{官方课件第 10 页；BIG-Bench}

IPA 例子展示了一个低频输出空间：输入是普通英语句子，目标却包含模型很少被要求按规则生成的音标序列。读图时应把“能否识别任务意图”“是否掌握音素映射”“能否稳定输出特殊字符”分开；任务准确率跨阈值只说明整条链终于可用，不等价于每个子能力同时诞生。

\subsection{Inverse scaling 为什么可能只是曲线前半段}

如果只看到 small model 优于 medium model，很容易得出“继续扩大会越来越差”的结论。课堂用一个 deliberately provocative 的词选择例子说明：模型变大后可能更忠实地复现训练分布中的常见关联，从而先偏离任务目标；再继续扩大后，更强的 instruction/context understanding 又可能把趋势拉回。

\lecturefigure{slide-11.jpg}{Inverse scaling can become U-shaped}{官方课件第 11 页；Wei, Tay, and Le, 2022}

\begin{knowledgebox}{读图：U-shape 的三段解释}
图中的 small model、medium model 与 large model 不应只被贴上“好/坏”标签。第一段可能是模型尚未学到强 prior；第二段是 prior 变强但任务控制不足；第三段是 context sensitivity、instruction following 或更高层组合能力追上 prior。真正的教学结论是：在曲线没有覆盖更大规模前，不要把局部负斜率外推成永久规律。
\end{knowledgebox}

\begin{warningbox}{稀疏 checkpoint 会制造视觉戏剧性}
如果横轴只有三四个模型，U-shape 的最低点和“阈值”都只能粗略定位。不同 prompt、seed、metric 或数据切分也可能移动曲线。正确表述应是“在被测试的模型族和设置中观察到趋势反转”，而不是“所有 inverse scaling 最终都会被规模解决”。
\end{warningbox}

\subsection{Emergent prompting technique：方法本身也有门槛}

前面讨论的是 task ability，课件随后把视角转向 intervention：同一个 RLHF 或 prompting 方法，在小模型上可能没有正收益，在大模型上才显示 delta。方法是否有效因此不能只在便宜的小模型上验证后直接外推。

\lecturefigure{slide-12.jpg}{一种干预可在小模型上有负收益、在大模型上转为正收益}{官方课件第 12 页；Wei et al., 2022}

\begin{knowledgebox}{读图：比较的是同规模内的 delta}
先在每个规模上比较 baseline 与 intervention，而不是只比较最左和最右的绝对分数。小模型处橙线低于蓝线，说明干预占用了模型尚不具备的表示或优化能力；大模型处橙线超过蓝线，说明相同干预终于能利用更强 base capability。这个例子也提醒我们区分 pretraining scale、prompting 与 weight-changing post-training。
\end{knowledgebox}

\teachervoice{学生问这些模型是否都经过 instruction fine-tuning，讲者明确回答这组 base-model emergence 图没有做 instruction fine-tuning。随后他用 RLHF 说明：只在小模型上试验一个方法并观察失败，并不足以证明它在更大模型上也无效。}

\subsection{本章小结}

Few-shot emergence 是相对 chance baseline 的任务级现象；inverse scaling 的局部负趋势可能在更大规模反转；prompting 或 post-training intervention 也可能有自己的能力门槛。三者共同反对一种常见错误：用有限小模型实验对整个规模区间作线性外推。

